\chapter{Additional features}
\label{chap_add_feat}
%%%%%%%%%%%%%%%%%%%%%%%%%%%%%%%%%%%%%%%%%%%%%%%%%%%%%%%%%%%%%%%%%%%%
\section{Data forwarding}

\section{Dynamic branch prediction}
\label{dyn_br}
In order not to waste too much clock cycles waiting for the correct value of the \textit{PC} register, we have designed a unit that provides a dynamic prediction regarding a branch instruction. In particular, it is possible, with a defined value of the \textit{PC} corresponding to a branch, to get whether is a better speculation to take or not to take the branch, while waiting for the correct \textit{PC}, which comes after the execution stage. 
There are actually 16 possible entries in our branch. Each of them is composed by a part of the \pc and a prediction state, represented on two bits, as reported in table \ref{pred_tab}. Every time a new instruction is issued, this entity checks whether it is a branch type or not. If so, the bits from 5 to 2 of the \pc related to the instruction are taken into account as address to the table, as shown in figure \ref{bht_fig}. At this point, the content of this location is checked:
\begin{itemize}
	\item if the value stored corresponds to the \pc, then the prediction contained is the one related to thie current instruction, and therefore it can be freely used
	\item if not, the previously stored value is overwritten, as long as its prediction: the current prediction is set to a "weakly not taken case"
\end{itemize}

\begin{table}[]
	\centering
	\begin{tabular}{l|cccc}
		\toprule
		& 00        & 01      & 10     & 11       \\
		\toprule
		Encoding     & Strong NT & Weak NT & Weak T & Strong T \\
		\midrule
		NS when HIT  & 00        & 00      & 11     & 11       \\
		NS when MISS & 01        & 10      & 01     & 10      	\\
		\bottomrule
	\end{tabular}
\caption{Possible states for predictor}
\label{pred_tab}
\end{table}


\begin{figure}
	\centering
	\includegraphics[scale=0.5]{chapters/figures/bht}
	\caption{Schematic for BHT}
	\label{bht_fig}
\end{figure}